\section{OS Foundation}
\label{sec:foundation}

Every operating system needs a foundation: a kernel that talks directly to
hardware, an init system that starts services at boot, a package manager for
installing software, and a base set of userspace tools. Building all of this
from scratch would take decades and produce something inferior to what already
exists. This project therefore builds on top of an existing Linux distribution
that provides the entire foundation layer, and adds everything above it: the
desktop environment, event system, knowledge graph, and AI layer.

\subsection{What is a Linux Distribution}
\label{sec:what-is-distro}

The \emph{Linux kernel} is the core of the operating system: the program that
manages hardware access, memory, processes, and the filesystem. It is developed
by Linus Torvalds and thousands of contributors worldwide and is released under
the GPL~v2 license. The kernel alone, however, is not a usable system. A
\emph{Linux distribution} bundles the kernel together with an init system, a
package manager, a C standard library, and a collection of userspace tools into
a coherent, installable operating system.

Different distributions make different choices about package freshness, update
cadence, governance, and default software selection. These choices have real
consequences for a project building on top: a distribution with stale packages
cannot support modern hardware or toolchains, and one with poor governance can
change direction in ways that break downstream projects without warning.

\subsection{Base Distribution}
\label{sec:base-distro}

This project builds on \emph{Fedora Linux}~\cite{fedora} as its base
distribution. Fedora is a community distribution backed by Red Hat
infrastructure but governed independently through the Fedora Council, which
makes binding decisions on project direction without requiring Red Hat
sign-off. Packages track upstream closely and a new major release ships every
six months. Btrfs has been the default filesystem since Fedora~33 (2020), and
the Flatpak and Wayland ecosystems receive stronger first-class support here
than on most other distributions.

The selection criteria were deliberately narrow: packages current enough for
modern toolchains and recent hardware without arriving unvalidated daily;
governance not concentrated in a single commercial entity; Btrfs as a
first-class filesystem; and broad enough hardware coverage to work on consumer
machines without manual driver installation. Fedora satisfies all four.

\emph{RPM} (RPM Package Manager) is the native package format. It is one of
the two dominant Linux package formats alongside Debian's \texttt{.deb}. A
package is a compressed archive containing the files to install, metadata about
name and version, and dependency declarations that the package manager resolves
automatically. The \emph{Open Build Service}~\cite{obs} (OBS), described in
Section~\ref{sec:infra}, handles building and signing RPMs for this project and
is fully compatible with Fedora as a build target.

The six-month release cadence means periodic major upgrades are required.
This is managed through the update strategy in Section~\ref{sec:updates}:
snapshots are taken before every upgrade, and the health check and rollback
mechanism make a failed major upgrade recoverable without user intervention.
Each Fedora release is supported for approximately thirteen months, giving a
comfortable window for validation before an upgrade is required.

\subsection{Update Strategy}
\label{sec:updates}

Updates on most operating systems involve some degree of risk: something breaks,
the user does not know why, and getting back to a working state requires either
expertise or luck. This system addresses that with a layered approach that makes
the worst-case outcome a single keypress rather than an afternoon in recovery
mode.

The approach is a mutable filesystem with automatic rollback protection.
\emph{Mutable} means the filesystem works the way Linux filesystems always have:
software can be installed freely, system files can be modified, and standard
developer tools work without restrictions. This distinguishes it from
\emph{immutable} distributions such as Fedora Silverblue and NixOS, where the
root filesystem is read-only and changes require going through a specific
toolchain. Immutable systems have genuine safety benefits, but they break
workflows that the target audience relies on: installing language runtimes with
\texttt{pip} or \texttt{cargo}, building projects from source, or running a
local development server. The safety benefits of immutability are achievable
through snapshot tooling without those costs.

Figure~\ref{fig:update-flow} shows the full update flow.

\begin{figure}[htbp]
\centering
\begin{tikzpicture}[
  block/.style={
    draw, thick, rectangle, rounded corners=3pt,
    minimum width=5cm, minimum height=0.85cm,
    font=\small, align=center, inner sep=5pt
  },
  decision/.style={
    draw, thick, diamond, aspect=3,
    minimum width=5cm, minimum height=0.85cm,
    font=\small, align=center, inner sep=4pt
  },
  outcome/.style={
    draw, thick, rectangle, rounded corners=3pt,
    minimum width=3.8cm, minimum height=0.85cm,
    font=\small, align=center, inner sep=5pt
  },
  arr/.style={->, thick},
]
\def\sep{1.15}
\node[block]    (avail) at (0, 4*\sep) {Update available};
\node[block]    (snap)  at (0, 3*\sep) {Pre-update snapshot (Snapper)};
\node[block]    (apply) at (0, 2*\sep) {Update applied};
\node[block]    (boot)  at (0, \sep)   {First boot after update};
\node[decision] (check) at (0, 0)      {Health check};
\node[outcome]  (ok)    at (-3.2, -2*\sep) {Continue normally};
\node[outcome]  (fail)  at ( 3.2, -2*\sep) {Auto-rollback + notify};
\draw[arr] (avail) -- (snap);
\draw[arr] (snap)  -- (apply);
\draw[arr] (apply) -- (boot);
\draw[arr] (boot)  -- (check);
\draw[arr] (check.west) -| node[pos=0.2, above, font=\footnotesize] {pass} (ok);
\draw[arr] (check.east) -| node[pos=0.2, above, font=\footnotesize] {fail} (fail);
\end{tikzpicture}
\caption{Update flow. A Snapper snapshot is created before every update.
If the post-update health check fails, the system rolls back to the snapshot
automatically and notifies the user.}
\label{fig:update-flow}
\end{figure}

\paragraph{Btrfs and copy-on-write.}
Most widely used Linux filesystems, including \emph{ext4}, use
\emph{overwrite-in-place} semantics: when a file is modified, the old data is
replaced at the same disk location. Once overwritten, the previous version is
gone.

\emph{Btrfs}~\cite{archwiki:btrfs}
(B-Tree Filesystem) uses \emph{copy-on-write} (CoW) semantics instead. When a
file is modified, Btrfs writes the new data to a different location on disk and
updates the metadata pointer. The old version remains intact until space is
explicitly reclaimed. This has two important consequences: filesystem snapshots
become nearly free in terms of disk space (a snapshot is a set of metadata
pointers to existing blocks, not a copy of the data), and the filesystem remains
self-consistent after a crash because incomplete writes never overwrite existing
valid data.

Btrfs was merged into the Linux kernel in 2009 and has been the default
filesystem on Fedora since 2020. Its CoW architecture is the technical
foundation for this project's rollback strategy.

\paragraph{Btrfs subvolumes.}
A \emph{subvolume} is a named, independently snapshotable region of a Btrfs
filesystem. It behaves like a directory but can be mounted separately, snapshotted
independently, and assigned its own quota limits. This project uses the following
layout, consistent with the convention established by Fedora:

\begin{lstlisting}[language=bash]
@          ->  mounted at /            (system root)
@home      ->  mounted at /home        (user data)
@snapshots ->  mounted at /.snapshots
@log       ->  mounted at /var/log
@cache     ->  mounted at /var/cache
\end{lstlisting}

Keeping \texttt{/} and \texttt{/home} on separate subvolumes means a system
rollback (restoring \texttt{@} to a previous snapshot) never touches user data
in \texttt{@home}. Log and cache data are excluded from system snapshots,
preventing large log files from inflating snapshot sizes.

\paragraph{Snapper.}
\emph{Snapper}~\cite{archwiki:snapper}
is the snapshot management tool that sits on top of Btrfs. It handles creating,
listing, comparing, and restoring snapshots automatically. Snapper is configured
per subvolume: the \texttt{@} subvolume gets automatic pre/post snapshots around
package manager operations and a timeline policy for background snapshots. The
\texttt{@home} subvolume has its own independent daily snapshot schedule.

Before any update is applied, Snapper creates a tagged pre-update snapshot of
the \texttt{@} subvolume. This is a precondition for the update, not an optional
step. A matching post-update snapshot is created after the package manager
finishes, forming a before/after pair that Snapper can diff to show exactly what
changed.

\paragraph{Bootloader integration.}
The bootloader is configured to list available snapshots as boot entries. If an
update produces a system that fails to boot, the firmware presents the previous
snapshot as a selectable entry on the next start. No terminal or recovery mode
is required. On Fedora this is handled by \texttt{grub2-snapper-plugin}.

\paragraph{Post-update health check.}
On the first boot after an update, a systemd oneshot service runs a health check
before the desktop session starts. It verifies that the display system
initializes, that the network stack is reachable, that the Graph Daemon and
Event Bus start successfully, and that no systemd services have entered a failed
state. If any check fails, the service triggers an automatic rollback by setting
the pre-update snapshot as the default boot entry and rebooting. The user
receives a plain-language notification on the rolled-back system.

\paragraph{Snapshot retention.}
Snapper's retention policy for this project removes snapshots older than seven
days while always retaining the five most recent. Both values are user-adjustable
in Settings. The \texttt{@home} subvolume retains the three most recent daily
snapshots; home snapshots are a safety net for accidental file deletion, not a
system recovery mechanism.

\subsection{System Infrastructure}
\label{sec:infra}

\paragraph{systemd.}
\emph{systemd}~\cite{archwiki:systemd}
is the init system and service manager on Fedora and on
most modern Linux distributions. When the kernel finishes booting, the first
userspace process it starts is \texttt{systemd} (PID~1), which reads
\emph{unit files} describing each service, its dependencies, and its startup
environment, and starts them in dependency order. systemd also handles logging
via the \emph{journal}, socket activation, timer-based scheduling, and hardware
event management.

All custom daemons in this project are deployed as systemd service units.
systemd's unit file format provides sandboxing options that are enforced at the
kernel level without any code changes in the service: \texttt{PrivateNetwork=true}
to deny all network access, \texttt{PrivateTmp} for an isolated temporary
directory, \texttt{MemoryDenyWriteExecute} to prevent writable executable memory,
and \texttt{SystemCallFilter} to allowlist only the syscalls the service
actually needs.

\paragraph{systemd-homed.}
\emph{systemd-homed}~\cite{archwiki:homed} is a systemd component
that manages user home directories as self-contained, portable, encrypted
images. Unlike traditional Unix home directories, which are plain directories
under \texttt{/home} identified by a numeric user ID, a systemd-homed home is a
\emph{LUKS-encrypted loopback image}: a file on disk that is mounted at login
(after the encryption key is derived from the user's credentials) and unmounted
at logout. When a user is not logged in, their home directory is an opaque
encrypted file; even root cannot read its contents without the user's passphrase.

This is the technical foundation for the profile isolation system described in
Section~\ref{sec:platform}: each profile is a Linux system user with its own
systemd-homed image. When a profile is not active, its image is unmounted and
its key is discarded from memory.

\emph{LUKS} (Linux Unified Key
Setup~\cite{archwiki:luks}) is
the standard disk encryption layer on Linux, built on top of the
\texttt{dm-crypt} kernel module. It provides a standardized header format for
encrypted block devices, support for multiple key slots so that several
passphrases or key files can unlock the same volume, and memory-hard key
derivation functions that make offline brute-force attacks expensive. LUKS is
what systemd-homed uses internally to encrypt each home image.

\paragraph{Session manager and greeter.}
User sessions are managed by \emph{greetd}~\cite{greetd}, a minimal login
daemon that handles authentication and session startup. greetd is
intentionally small: it authenticates the user via PAM, launches the
configured session command on success, and exits. All policy decisions, visual
presentation, and user-facing interaction are delegated to a separate
\emph{greeter} process that communicates with greetd over a Unix socket using a
simple JSON protocol.

This clean separation between authentication daemon and login UI is why greetd
was chosen over display managers such as GDM or SDDM, which bundle both
concerns into a single process with their own GUI frameworks. The login screen
in this project is a first-party Tauri application that implements the greetd
IPC protocol: it presents the user interface, collects credentials, and passes
them to greetd for PAM verification. The Tauri greeter has access to nothing
other than the greetd socket; it runs before any user session is active, with
no access to the graph, the event bus, or any user data.

\paragraph{Package management.}
At the base level, \texttt{dnf}~\cite{dnf} handles dependency resolution,
cryptographic package signature verification, and repository management.

First-party project packages are built and distributed through the \emph{Open
Build Service} (OBS~\cite{obs}). OBS is a hosted build infrastructure
originally developed by SUSE and used by many open-source projects across
distribution families. Developers push package source files to OBS; OBS builds
the resulting RPMs in a clean Fedora chroot for every configured release target
and architecture, signs them with the project GPG key, and hosts the resulting
repository. This keeps the build pipeline reproducible and independent of the
developer's local environment.

\paragraph{Flatpak.}
Store applications are distributed as
\emph{Flatpak}~\cite{archwiki:flatpak}
packages. Flatpak is a distribution-agnostic packaging format for desktop
applications. A Flatpak bundles an application together with its specific
library dependencies into a sandboxed unit that runs isolated from the rest of
the system using Linux namespaces and \emph{bubblewrap}, with controlled access
to host resources via the \emph{XDG Desktop Portal}.

The alternative formats are \emph{Snap}, which routes all distribution through
Canonical's servers and has faced criticism for this centralization, and
\emph{AppImage}, which is a self-contained executable with no sandboxing and no
integrated update mechanism. Flatpak was chosen for its sandboxing model,
distribution neutrality, and strong support in both Fedora and the broader Linux
desktop ecosystem.

System-level components ship as native RPMs because they require tight
integration with the system: daemons must run outside any sandbox, and systemd
unit files must be installed to system paths. The Store handles both formats
transparently.

\subsection{Hardware Requirements}
\label{sec:hardware}

Three hard requirements determine the minimum supported hardware. They are not
adjustable without significant changes to the security and graphics stack and
are treated as hard constraints. In practice they exclude hardware older than
roughly 2018--2019.

\begin{table}[htbp]
\caption{Minimum hardware requirements}
\label{tab:hardware}
\begin{tabularx}{\linewidth}{lX}
\toprule
\textbf{Requirement} & \textbf{Reason} \\
\midrule
Linux kernel 6.6 or newer &
  Required for the eBPF program types used by the event monitoring system
  (Section~\ref{sec:ebpf}): \texttt{BPF\_PROG\_TYPE\_TRACEPOINT} for syscall
  attachment, \texttt{BPF\_PROG\_TYPE\_PERF\_EVENT} for process lifecycle,
  and CO-RE (Compile Once, Run Everywhere) via BTF type information for
  portable eBPF binaries. Kernel 6.6 is the LTS release that ships with
  current Fedora. \\
Vulkan-capable GPU &
  Required for hardware-accelerated compositing and for Windows compatibility
  via DXVK and vkd3d-proton. Covered automatically on AMD and Intel via Mesa;
  Nvidia requires the proprietary driver, with installer guidance provided. \\
systemd 255 or newer &
  Required for full systemd-homed portable home image functionality used by
  the profile system. \\
\bottomrule
\end{tabularx}
\end{table}

\paragraph{eBPF program types and CO-RE.}
The event monitoring system uses three eBPF program types. Tracepoint programs
attach to kernel syscall entry and exit points (\texttt{openat},
\texttt{execve}, \texttt{connect}, and others) and fire on every matching
call across all processes. Perf event programs handle process lifecycle events
including fork and exit. Ring buffer maps transfer collected events to
userspace with low overhead and without copying data through the kernel I/O
path.

All eBPF programs are compiled to eBPF bytecode at build time using the
\texttt{aya}~\cite{aya} framework and ship as pre-compiled binaries in the
system image. No compiler is required on the end-user's machine. The programs
use \emph{CO-RE} (Compile Once, Run Everywhere)~\cite{ebpf:core}, a mechanism
introduced in Linux 5.2 that uses BTF (BPF Type Format) kernel debug
information to make eBPF programs portable across kernel versions without
recompilation. Kernel 6.6 guarantees that the required BTF information is
present and that the specific tracepoint hooks this system attaches to are
stable. Kernels older than 5.8 lack ring buffer map support
entirely~\cite{ebpf:ringbuf}; kernels between 5.8 and 6.5 may work but are
not tested and not supported.

\paragraph{Vulkan.}
\emph{Vulkan}~\cite{vulkan} is a low-level graphics and
compute API by the Khronos Group, designed as the modern replacement for OpenGL.
It provides explicit control over the GPU at the cost of a more verbose API.
This project requires Vulkan for two reasons: the compositor uses Vulkan
for hardware-accelerated compositing, and the Windows compatibility layer
(Section~\ref{sec:gaming}) translates Direct3D calls to Vulkan via DXVK and
vkd3d-proton. On AMD and Intel hardware, Vulkan support is provided by
\emph{Mesa}~\cite{mesa}, the open-source graphics
driver stack included with Fedora. Nvidia hardware requires
Nvidia's proprietary driver because the open-source \texttt{nouveau} driver does
not provide Vulkan support on most hardware generations.

\paragraph{Hardware support scope.}
The design principle ``works on first boot'' requires an honest statement of
scope. It means: on supported hardware classes, the system boots, the desktop
starts, networking works, and audio works without manual configuration.

The supported hardware classes where this principle holds are mainstream x86-64
laptops and desktops with Intel or AMD CPUs released after 2018, with
integrated or discrete Intel or AMD graphics. This covers the large majority
of consumer hardware available today.

Two known categories require additional handling. Nvidia discrete GPUs require
the proprietary driver~\cite{fedora:nvidia}. RPM Fusion is configured as a
repository source in the system image; on first boot, the installer detects
Nvidia hardware and installs the proprietary driver automatically before the
user reaches the desktop. No manual repository configuration is required.
Broadcom Wi-Fi adapters, found in some Apple hardware and a smaller set of
laptops, require a proprietary driver that is not included in the Fedora
kernel~\cite{fedora:broadcom}. The same first-boot detection logic applies
where a driver package is available; hardware without a packaged driver is
documented explicitly in the compatibility list rather than discovered by
users after installation.

\subsection{Configuration Format}
\label{sec:config}

All system components and first-party applications use
\emph{TOML}~\cite{toml} (Tom's Obvious Minimal Language) as
the configuration format. TOML was designed by Tom Preston-Werner (co-founder of
GitHub) to be unambiguous and easy to edit without special tooling.

The two most common alternatives are \emph{YAML} and \emph{JSON}. YAML is
human-readable and widely used, but its indentation-sensitive syntax is a
consistent source of subtle errors: a single misaligned space changes the meaning
of the document, and certain values parse unexpectedly (bare \texttt{NO} is a boolean in YAML~1.1, not a string; YAML~1.2 fixes this, but library support for 1.2 is inconsistent in practice). JSON is unambiguous but has no comment
support and is more verbose than necessary for human-edited files.

TOML has clear, explicit syntax with no indentation sensitivity, supports inline
comments, and has an unambiguous type system: strings, integers, floats,
booleans, datetimes, arrays, and tables. Using one format consistently means a
developer who reads one configuration file can read all of them, and validation
and migration tooling only needs to be written once.

Personal settings live at
\texttt{\textasciitilde/.config/<app-id>/config.toml} and system-wide settings
at \texttt{/etc/<service-name>/config.toml}, following the
\emph{XDG Base Directory Specification}~\cite{xdg:basedir}. Some third-party applications expect other formats; a
small helper translates without moving the source of truth:

\begin{lstlisting}[language=bash]
# Export as environment variables for a legacy app
eval $(os-config-helper export \
    --toml ~/.config/myapp/config.toml --format env)

# Export as JSON for tools that require it
os-config-helper export \
    --toml ~/.config/myapp/config.toml --format json
\end{lstlisting}

The TOML file is always the authoritative source; the helper is a read-only
translation layer with no business logic of its own.
